\section{Complex Data: Sets, Curves, and Ensembles}

Classic multidimensional data has records with numerical and categorical
attributes. \term{Complex data} adds richer attribute kinds: \term{sets},
\term{curves} (function graphs / time series), and \term{surfaces}. This
section covers three such cases that recur in the exam: set-typed data
(how to visualise it, with UpSet as the modern go-to answer), families of
function graphs (Konyha et al.), and simulation ensembles (Matkovic et al.).
All three are united by one idea: an attribute is no longer a single scalar
but a structured object, and the same visual-analysis toolkit
(\term{coordinated multiple views} + \term{brushing and linking}) is used to
explore them.

% ======================================================================
\subsection{Set-Typed Data}

\begin{keybox}[Definition]
\term{Set-typed data}: a collection of \term{elements}, each of which has a
membership in one or more of several \emph{possibly overlapping} \term{sets}.
Equivalently, each element carries a \term{set-typed attribute} whose value is
a subset of the universe of sets (a \term{set family} / set system over the
elements). Sets may relate by \term{containment, exclusion, and intersection}.
\end{keybox}

\textbf{Example (movies).} Universe of sets $=$ genres
$\{\text{Romance},\text{Drama},\text{Comedy},\dots\}$. Each movie is an element;
its set-typed attribute is its genre set, e.g.
\textit{Cinema Paradiso} $=\{$Romance, Drama, Comedy$\}$, \textit{Touch}
$=\{$Romance$\}$. Other real-world cases: club memberships, product features,
employee skill sets, search-engine result overlaps, patient symptoms.

\textbf{Key notions} (Sets-STAR, Alsallakh et al.).
\begin{itemize}
\item \term{Set membership degree} of an element $=$ number of sets it
  belongs to (its \term{degree}).
\item \term{Exclusive membership}: an element belongs to exactly one set / one
  specific intersection.
\item \term{Intersection / overlap}: the set of elements shared by a group of
  sets. With $n$ sets there are up to $2^n$ possible membership combinations,
  which is why naive approaches do not scale.
\item Data representations: cardinalities only, multi-valued attribute
  (adjacency-list style), membership lists (edge-list style), or Boolean
  attributes (adjacency-matrix style).
\end{itemize}

\textbf{Typical tasks} (what a technique must support): find elements of a set,
find sets of an element, filter by membership degree; find the number of sets;
analyse inclusion / exclusion / intersection relations; identify the sets in a
given intersection and the largest / smallest intersections; compare set and
intersection cardinalities and similarities.

\subsubsection{Techniques (Sets-STAR categories)}

\begin{algobox}[Euler / Venn diagrams]
\textbf{Idea:} sets $=$ labelled closed curves (circles, ellipses, polygons);
set relations $=$ curve overlaps. A \term{Venn} diagram shows \emph{all} $2^n$
possible overlaps; an \term{Euler} diagram shows only the overlaps that
actually occur (it can be \term{well-matched}: regions reflect the true
relations).\\
\textbf{Reads:} membership pops out preattentively via enclosure.\\
\textbf{Trade-offs:} \emph{does not scale beyond $\sim$3 sets}; for $\ge 4$
sets, well-matched + well-formed (simple, smooth, area-proportional) curves are
often impossible, and \emph{some intersections become undrawable}. Area-judgement
bias hurts area-proportional variants.
\end{algobox}

\begin{algobox}[Over-plotting / point-based (overlays)]
\textbf{Idea:} elements live in an existing layout (scatterplot, map, graph)
and set membership is overlaid with \term{colour/glyphs} (coloured dots,
pie-glyphs, hatching) or with \term{region} overlays (\term{Bubble Sets},
texture splatting) or \term{line} overlays (\term{LineSets}, \term{Kelp
Diagrams}).\\
\textbf{Use when:} elements have a meaningful spatial reference and set
membership is secondary. Region overlays handle $\sim$4--20 sets.\\
\textbf{Trade-offs:} limited number of sets/elements; layout artifacts
(spurious overlaps, crossings); overlays cannot move the elements.
\end{algobox}

\begin{algobox}[Matrix-based]
\textbf{Idea:} a \term{set-membership matrix} -- rows $=$ elements, columns $=$
sets (cells mark membership); or sets $\times$ sets cells holding a similarity /
overlap measure (heatmap). Reordering rows/columns reveals clusters of similar
membership. Examples: ConSet, PixelLayers, frequency grids.\\
\textbf{Trade-offs:} fairly scalable in both elements and sets, clutter-free
(no edge crossings), but limited in which set \emph{relations} they express,
and patterns are \emph{sensitive to ordering}.
\end{algobox}

\begin{algobox}[Node-link / graph-based]
\textbf{Idea:} model element--set membership as a \term{bipartite graph} (edges
$=$ memberships). Examples: Jigsaw lists, anchored maps, PivotPaths; Radial
Sets / Circos for set--set links.\\
\textbf{Trade-offs:} elements are visible as individual objects and clusters of
similar membership show up, but \emph{limited scalability due to edge
crossings}; element--set diagrams show no set relations directly.
\end{algobox}

\begin{algobox}[Aggregation-based]
\textbf{Idea:} when elements are many, follow ``overview first'': aggregate
elements into a single visual element encoding a \emph{frequency}. Bar charts of
set/intersection sizes, \term{Set'o'gram} (Freiler et al.), Mosaic / Double-Decker
plots, Parallel Sets, \term{Radial Sets} (Alsallakh et al.).\\
\textbf{Trade-offs:} \emph{highly scalable in number of elements}; some show how
attributes correlate with membership; but usually do not show single elements
as objects and are limited in set relations.
\end{algobox}

\begin{keybox}[UpSet -- the modern go-to technique]
Lex, Gehlenborg, Strobelt, Vuillemot \& Pfister, \emph{UpSet: Visualization of
Intersecting Sets}, IEEE TVCG 2014. \textbf{This is the standard answer to
``how do you visualise set-typed data with many sets?''}
\end{keybox}

UpSet replaces the overlapping regions of an Euler/Venn diagram with a
\term{matrix of intersections} plus \term{bar charts}:
\begin{itemize}
\item \textbf{Set-size bars} -- one bar per set giving each set's cardinality.
\item \textbf{Intersection matrix} -- columns are the (non-empty) intersections;
  each column is a vertical strip of dots over the set rows. A \term{filled
  dot} $=$ ``in this set'', an \term{empty dot} $=$ ``not in this set''; dots in
  the same column are joined by a line. A column with one filled dot $=$
  ``only in that set''; all dots filled $=$ ``in all sets''.
\item \textbf{Intersection-size bars} -- a bar chart giving the cardinality of
  each intersection, usually \emph{sorted} (e.g. by size) so the dominant
  combinations are obvious.
\end{itemize}
Because intersections are laid out as a sortable list rather than packed into
the plane, UpSet \emph{scales to many sets and many elements} and makes the
hard set tasks (which sets are in an intersection? which intersection is
largest?) directly readable. Related slide techniques: OnSet (binary
set matrices with set operations), Radial Sets, Set'o'gram.

\begin{pitfall}[Common deductions]
Do not draw a Venn diagram for $\ge 4$ sets and claim it ``works'' -- some
intersections are undrawable and area judgements are biased. Do not show set
sizes as a \term{pie chart}: set categories are not mutually exclusive and do
not sum to a whole. For ``many sets'' always reach for \textbf{UpSet} (or
aggregation/matrix methods), not Euler/Venn.
\end{pitfall}

\begin{center}
\begin{tikzpicture}[x=1cm,y=1cm,font=\scriptsize]
  % --- geometry: matrix rows A,B,C at y=2,1,0 ; columns x=0..4 ---
  \def\rA{2} \def\rB{1} \def\rC{0}
  % intersection-size bars (top), sorted descending
  % columns and the heights of their size bars
  \foreach \x/\h in {0/1.7,1/1.3,2/1.0,3/0.7,4/0.45}{
    \fill[tuwblue!70] (\x-0.18,2.55) rectangle (\x+0.18,{2.55+\h});}
  \node[anchor=south,rotate=90,tuwdark] at (-1.05,3.4)
        {\,intersection size};
  % --- matrix dots ---
  % membership pattern per column (filled rows). 1=filled
  % col0: A   col1: A,B   col2: B,C   col3: A,B,C   col4: C
  \foreach \x in {0,1,2,3,4}{
    % empty dots everywhere first
    \foreach \r in {\rA,\rB,\rC}{
      \draw[gray!55,fill=gray!12] (\x,\r) circle (0.16);}}
  % filled dots + connecting lines per column
  % col0: A
  \fill[tuwdark] (0,\rA) circle (0.16);
  % col1: A,B
  \draw[tuwdark,line width=1pt] (1,\rB)--(1,\rA);
  \fill[tuwdark] (1,\rA) circle (0.16); \fill[tuwdark] (1,\rB) circle (0.16);
  % col2: B,C
  \draw[tuwdark,line width=1pt] (2,\rC)--(2,\rB);
  \fill[tuwdark] (2,\rB) circle (0.16); \fill[tuwdark] (2,\rC) circle (0.16);
  % col3: A,B,C
  \draw[tuwdark,line width=1pt] (3,\rC)--(3,\rA);
  \fill[tuwdark] (3,\rA) circle (0.16); \fill[tuwdark] (3,\rB) circle (0.16);
  \fill[tuwdark] (3,\rC) circle (0.16);
  % col4: C
  \fill[tuwdark] (4,\rC) circle (0.16);
  % --- row labels ---
  \node[anchor=east] at (-0.55,\rA) {$A$};
  \node[anchor=east] at (-0.55,\rB) {$B$};
  \node[anchor=east] at (-0.55,\rC) {$C$};
  % --- set-size bars (left, horizontal) ---
  \foreach \r/\w in {\rA/1.3,\rB/1.6,\rC/1.0}{
    \fill[keygreen!70] (-2.4,\r-0.16) rectangle (-2.4+\w,\r+0.16);}
  \node[anchor=south,tuwdark] at (-1.9,2.5) {set size};
\end{tikzpicture}
\end{center}
\sketchnote{An UpSet plot has three coupled parts. \emph{Top-left:} a small bar
chart of total set sizes (sets $A,B,C$). \emph{Bottom-left:} the set rows
$A,B,C$ as labels. \emph{Centre/bottom:} a dot matrix -- one column per
intersection, each column a vertical run of filled (member) / empty (non-member)
dots, filled dots in a column joined by a line. \emph{Top:} a horizontal bar
above each column whose length is that intersection's cardinality, columns
sorted by size. Reading a column's dots tells you exactly which
sets define the intersection; the bar tells you how big it is.}

\subsubsection{Comparison of set-visualisation categories}

\begin{center}
\begin{tabular}{@{}p{2.6cm}p{1.9cm}p{2.1cm}p{2.4cm}p{3.0cm}@{}}
\toprule
\textbf{Category} & \textbf{Scales in \#sets} & \textbf{Scales in \#elements}
  & \textbf{Shows intersections?} & \textbf{Order matters?}\\
\midrule
Euler / Venn & $\sim$3 (poor) & tens & yes, but some undrawable for $\ge 4$
  & layout-driven\\
Overlays (point/region/line) & 4--20 & tens--hundreds & partially (regions
  show overlap) & element positions fixed by base view\\
Node-link (bipartite) & tens & hundreds & set--set only (Circos/Radial); not in
  element--set & sensitive (crossings)\\
Matrix-based & up to $\sim$100 & up to $\sim$100s & limited relations
  & \textbf{yes} (reordering reveals patterns)\\
Aggregation (Set'o'gram, Radial Sets, Double-Decker) & tens--$\sim$50 & large
  (hundreds+) & limited / on demand & yes\\
\textbf{UpSet} & \textbf{many} & \textbf{many} & \textbf{yes, explicitly +
  sizes} & \textbf{sortable} (by size/degree)\\
\bottomrule
\end{tabular}
\end{center}

% ======================================================================
\subsection{Families of Curves / Function Graphs}

\begin{keybox}[Definition]
A \term{function graph} $f(\mathbf{x},t)$ is a dependent variable plotted over
an independent variable (e.g. time): one 2D curve per data record. A
\term{family of function graphs} is the set of all such curves for one
dependent variable across all records -- many curves drawn in the same view.
Source: Konyha, Matkovic, Gra\v{c}anin, Jelovi\'c \& Hauser, \emph{Interactive
Visual Analysis of Families of Function Graphs}, IEEE TVCG 2006.
\end{keybox}

\textbf{Data model.} Records have $m$ \term{independent variables}
$\mathbf{x}=[x_1,\dots,x_m]$ and $n$ \term{dependent variables}. Dependents are
either \term{regular} (a single scalar per record) or \term{function-graph}
variables (a value for each time $t$, i.e. a curve). A curve can be treated
either as an \term{atomic item} or as an \term{attribute / dimension} of a
record (cf. the slides' table with a \texttt{Pressure(t)} / \texttt{Velocity(t)}
column). Examples: traffic-sensor occupancy/volume over a day (the Minneapolis
data set: $\sim$94{,}752 graphs, 144 values each), animal paths, ergonomics,
fuel-injection rate curves.

\textbf{Core problem: over-plotting.} Thousands of curves
(e.g. 45{,}000) overlap into a dense band, so individual curves and
regions of maxima/minima are hidden.

\begin{algobox}[Interactive visual analysis of curve families]
\textbf{Setup:} the \term{curve view} (a family of function graphs) embedded in
\term{multiple coordinated linked views} alongside histograms, scatterplots,
parallel coordinates, and a map of the independent variables. Brushing in any
view highlights the same items (\term{focus}) in all others; context is shown
desaturated (\term{focus+context}).\\
\textbf{Density rendering:} draw curves semi-transparently so pixels through
which more graphs pass are rendered with higher intensity -- a \term{banded /
density representation} that reveals the shape of a dense family.\\
\textbf{Brushing in function space:}
\begin{itemize}
\item \term{Line brush} -- a line segment in the curve view selects all curves
  that \emph{cross} it. Intuitive, easy to define, would be far more complex as
  an SQL query.
\item \term{Rectangular / polyline brush} -- selects curves passing through a
  region (or, with edge constraints, curves with a local max/min inside it).
\end{itemize}
\textbf{Brushing in derived attribute / parameter spaces:} aggregate each curve
into scalar \term{curve features} (max, min, slope, value-at-$x$, mean) and
brush those in scatterplots/histograms; link back to the curves and to the
independent variables.\\
\textbf{Composite brushing:} combine brushes with logical \term{AND, OR, SUB}
to iteratively drill down or broaden the focus.
\end{algobox}

\textbf{Analysis procedures.} \term{Black-box reconstruction}: understand how
dependent curves change as independent variables vary
(overview $\to$ zoom/filter $\to$ details). \term{Inverse analysis}: brush a
desired curve shape, then find which independent-variable combinations produce
it. \term{Multidimensional relations}: brush in one curve family and study the
linked curves of another. A \term{colour gradient} along a brush links brushed
items across views.

% ======================================================================
\subsection{Simulation Ensembles}

\begin{keybox}[Definition]
A \term{simulation ensemble} is a set of many \term{simulation runs}
(ensemble \term{members}), each produced by the same model with different
\term{control parameters / inputs}. Formally a member is a pair
$(\mathbf{x}^i,\mathbf{y}^i)$ where $\mathbf{x}\in\mathbb{R}^n$ are control
parameters and $\mathbf{y}=S(\mathbf{x})$ the outputs; the ensemble is
$E=\{(\mathbf{x}^i,\mathbf{y}^i)\}$. Members are \term{multi-dimensional and
often time-dependent} (scalars, curves, surfaces, fields). Source: Matkovic,
Gra\v{c}anin \& Hauser, \emph{Visual Analytics for Simulation Ensembles}, WSC 2018.
\end{keybox}

\textbf{Why ensembles.} Two questions drive simulation: (1) what is the output
for a given parametrisation (easy -- run it), and (2) which parametrisation
gives a desired output (hard -- inverting the model). Realistic models cannot be
inverted analytically, so we simulate \emph{many} parametrisations and explore
the ensemble. Sampling of the parameter space uses full-factorial or
low-discrepancy (Sobol) designs (e.g. the common-rail Diesel injector:
$5\times5\times5\times5\times7 = 4375$ runs).

\textbf{Goals / tasks.} \term{Parameter sensitivity} -- relate outputs to
control parameters and find which parameters matter; \term{model
reconstruction} -- find parameters for a desired output; \term{comparison} of
outputs from different parameter regions; identifying \term{representative} and
\term{outlier} members.

\begin{algobox}[Visual analytics for ensembles]
\textbf{Coordinated multiple views (CMV):} link \emph{parameter space} (left,
scatterplots/histograms/parallel coordinates of control parameters), the
\emph{curve/time views} (centre, e.g. needle-lift curves), and other
\emph{derived-output} views (right). A consistent left-to-right
parameter$\to$output layout helps build a mental model.\\
\textbf{Brushing \& linking across members:} brush a region of parameter space
or a desired output range; corresponding members highlight everywhere. Use the
curve \term{line brush} (cross-the-line selection) on the time-dependent
outputs.\\
\textbf{Aggregation:} reduce complex (non-scalar) outputs to \term{derived
scalar features} (min/max of a curve, etc.); support \term{on-the-fly}
aggregation during analysis so the session need not be restarted.\\
\textbf{Interactive ensemble steering:} start with a coarse sampling, then
\emph{initiate new simulation runs} from within the analysis where more detail
is needed (the ensemble grows as you explore). Three loops: inner analysis
loop (linking/brushing in CMV), ensemble-refinement loop (new parametrisations
/ rows), model-refinement loop (new model parameters / columns). \term{Hybrid
steering} adds an automatic regression-model optimisation to guide the
interactive search.\\
\textbf{Reproducibility:} structure the brushing space (snap-to-grid,
percentile grids), overlay statistics, and quantify findings so qualitative
insights (``the curves rise steeply'') become reproducible numbers.
\end{algobox}

\begin{center}
\begin{tikzpicture}[x=1cm,y=1cm,font=\scriptsize,>={Stealth[length=2mm]}]
  % ---- LEFT: parameter-space scatterplot with a brushed region ----
  \draw[rulec,thick] (0,0) rectangle (2.2,2.2);
  \node[anchor=south,tuwdark] at (1.1,2.25) {parameters};
  \fill[examorange!18] (0.9,0.7) rectangle (1.8,1.6); % brush
  \foreach \p in {(0.3,0.4),(0.6,1.7),(1.4,1.9),(0.4,1.2)}{
    \fill[gray!55] \p circle (1.3pt);}
  \foreach \p in {(1.1,1.0),(1.5,1.3),(1.0,1.4),(1.6,0.9),(1.3,1.2)}{
    \fill[examorange] \p circle (1.3pt);}
  % ---- CENTRE: curve view (output curves), brushed in orange ----
  \draw[rulec,thick] (3.0,0) rectangle (6.0,2.2);
  \node[anchor=south,tuwdark] at (4.5,2.25) {output curves};
  \draw[gray!50] plot[smooth] coordinates{(3.1,0.5)(3.9,1.3)(4.7,1.6)(5.9,1.7)};
  \draw[gray!50] plot[smooth] coordinates{(3.1,0.4)(3.9,0.8)(4.7,1.0)(5.9,1.05)};
  \draw[examorange,thick] plot[smooth]
        coordinates{(3.1,0.6)(3.9,1.7)(4.7,2.0)(5.9,1.95)};
  \draw[examorange,thick] plot[smooth]
        coordinates{(3.1,0.55)(3.9,1.55)(4.7,1.85)(5.9,1.8)};
  % ---- RIGHT: derived-output scatterplot ----
  \draw[rulec,thick] (6.8,0) rectangle (9.0,2.2);
  \node[anchor=south,tuwdark] at (7.9,2.25) {derived};
  \foreach \p in {(7.1,0.5),(7.4,1.8),(8.6,0.7),(7.2,1.1)}{
    \fill[gray!55] \p circle (1.3pt);}
  \foreach \p in {(8.0,1.5),(8.4,1.7),(8.1,1.2),(8.5,1.4)}{
    \fill[examorange] \p circle (1.3pt);}
  % ---- linking arrows ----
  \draw[->,tuwblue] (2.2,1.15) -- (3.0,1.15);
  \draw[->,tuwblue] (6.0,1.15) -- (6.8,1.15);
  % ---- data table below (brushed runs) ----
  \draw[rulec,thick] (3.0,-1.1) rectangle (6.0,-0.2);
  \foreach \y in {-0.45,-0.7,-0.95}{\draw[gray!40] (3.0,\y)--(6.0,\y);}
  \draw[gray!40] (4.0,-0.2)--(4.0,-1.1); \draw[gray!40] (5.0,-0.2)--(5.0,-1.1);
  \node[anchor=south,tuwdark] at (4.5,-0.18) {brushed runs (table)};
\end{tikzpicture}
\end{center}
\sketchnote{Ensemble CMV: left column $=$ parameter-space scatterplots/histograms
($dT_v, dT_p, P_{low},\dots$) with a brushed region; centre $=$ a curve view of
all members' output curves, brushed members in orange, context in grey; right
$=$ scatterplots of derived scalar outputs; a data table below lists the brushed
runs. Brushing any view re-highlights the same members everywhere.}

% ======================================================================
\subsection{Exam Angle}

\begin{exambox}[Set-typed data: strategy / list questions]
\textbf{``What is set-typed data and how would you visualise it?''} (1) Define:
elements with membership in multiple, possibly overlapping sets (set-typed
attribute). (2) List techniques by category: Euler/Venn, overlays
(point/region/line), matrix-based, node-link, aggregation. (3) State the
key trade-off: Euler/Venn fails beyond $\sim$3 sets (undrawable
intersections, area bias). (4) \textbf{Recommend UpSet} for many sets:
intersection matrix + sorted intersection-size bars + set-size bars; scales,
shows intersections explicitly. Mention the pie-chart pitfall and that pattern
visibility in matrices depends on ordering.
\end{exambox}

\begin{exambox}[Reading an UpSet plot -- visualisation literacy]
You may be shown an UpSet plot and asked to read it. Steps: (a) the left bars
give each \emph{set's} total size; (b) pick a column in the matrix -- the
\emph{filled} dots name the sets the elements belong to, \emph{empty} dots the
sets they do not; (c) the bar above that column gives the \emph{intersection's
size}. So ``only $B$'' $=$ one filled dot on $B$; ``$A\cap B$ but not $C$'' $=$
filled $A,B$, empty $C$; ``$A\cap B\cap C$'' $=$ all filled. The tallest
intersection bar is the dominant combination. Equivalent question for curves:
``a line brush selects which curves?'' -- all curves that \emph{cross the
line}.
\end{exambox}

\begin{exambox}[Curves \& ensembles]
Expect short-answer questions: define a \term{family of function graphs}
(many curves of one dependent variable; problem $=$ over-plotting; remedies $=$
density/banded rendering, line/composite brushing, brushing on derived curve
features in linked views). Define a \term{simulation ensemble} (many runs with
varied parameters; members multi-dimensional/time-dependent; goals $=$
parameter sensitivity, model reconstruction, representative/outlier members;
technique $=$ CMV linking parameter space $\leftrightarrow$ spatial/curve views,
aggregation, brushing across members, plus interactive ensemble steering).
\end{exambox}
